% Общие данные опубликованного прогона; это не самостоятельный документ.
% Источник: docs/simulation/ngspice-2026-11-28/report.json и run_info.json.
\textbf{Условия.} Прогон ngspice 45.2 от 28.11.2026 построен по XML KiCad:
20 Гц–20 кГц, 1003 отсчёта на канал, внешняя нагрузка 47 кОм.
Нормировка RIAA — на 1 кГц; эталон имеет постоянные времени
3180, 318 и 75 мкс. Источник сигнала идеальный дифференциальный.
\par\medskip
\textbf{Модель ОУ.} Линейное однополюсное приближение:
\[
 A(s)=\frac{A_0}{1+sA_0/(2\pi\,GBW)},\qquad
 A_0=10^5,\quad GBW=12\,\text{МГц}.
\]
В снимке KiCad сопротивления собраны до 23.51 и 90.91 кОм;
аналитическая конфигурация использует округлённые 23.5 и 90.9 кОм.
Параметры питания и выходы модели идеальны.
\par\medskip
\textbf{Результаты обоих каналов.}
\begin{center}
\begin{tabular}{@{}lr@{}}
Усиление на 1 кГц, В/В & 100.23888\\
Усиление на 1 кГц, дБ & 40.02072\\
Мин. ошибка RIAA (20 кГц), дБ & $-0.11942$\\
Макс. ошибка RIAA (382.36 Гц), дБ & $+0.00131$\\
\end{tabular}
\end{center}
При входе 4 мВ RMS между плечами выход на 1 кГц составляет
около 0.401 В RMS. AC-источники $+0.5$ и $-0.5$ В задают единичный
дифференциальный сигнал для расчёта передаточной функции;
это не проверка допустимого большого сигнала.
\par\medskip
\textbf{Как читать график.} Сверху — усиление с требуемым спадом RIAA;
в середине — фаза сигнала; снизу — отклонение формы АЧХ от эталона
после нормировки обеих характеристик на 1 кГц.
Кривые L/R совпадают из-за одинаковых номиналов и параметров модели,
без заданного разброса компонентов. Конечное усиление ОУ немного меняет
результат относительно идеального расчёта.
\par\medskip
\textbf{Границы.} Не учтены импеданс головки и кабеля, паразитные параметры
платы, шум, THD и ограничение выходного размаха. Совпадение каналов в модели
не задаёт реальное согласование или CMRR; запас устойчивости не проверен.
